\chapter{Introduction}\label{ch-intro}

\section{The spectral sum problem}\label{sec-problem}
Throughout this thesis a \emph{graph} means a finite, simple, undirected graph
without loops or multiple edges. For a graph $G$ of order $n$ we write
\[
 \lambda_1(G)\ge\lambda_2(G)\ge\cdots\ge\lambda_n(G)
\]
for the eigenvalues of its adjacency matrix $A_G$, listed in algebraically
decreasing order and repeated according to multiplicity. For an integer
$k$ with $1\le k\le n$ put
\[
 \Lambda_k(G)=\lambda_1(G)+\cdots+\lambda_k(G),
 \qquad
 M_k(n)=\max\{\Lambda_k(G):|V(G)|=n\},
\]
and, following Nikiforov~\cite{Nik15},
\[
 c_k=\sup\Bigl\{\frac{\lambda_k(G)}{|V(G)|}:|V(G)|\ge k\Bigr\},
 \qquad
 \tau_k=\limsup_{n\to\infty}\frac{M_k(n)}n .
\]
The signs in $\Lambda_k$ are retained: for $k\ge2$ the summands may be
negative, and $\Lambda_k$ is therefore not a norm. Two classical facts make
the extremal problem non-trivial and keep it anchored to $n$. First,
$\sum_i\lambda_i(G)=\tr A_G=0$ and $\sum_i\lambda_i(G)^2=2|E(G)|$, so
$\Lambda_k$ is bounded above by $\sqrt{k\cdot 2|E(G)|}\le\sqrt{k}\,n$ through
Cauchy--Schwarz; second, $\Lambda_n(G)=0$ identically. The interesting range
is thus $2\le k\ll n$, where the question posed by Hong~\cite{Hong93} asks
for the best possible order-linear coefficient.

The case $k=1$ is the spectral radius problem and is classical:
$\lambda_1(G)\le n-1$ with equality only for $K_n$. The case $k=2$ was
settled only recently~\cite{KLMPT,HW}. The case $k=3$ is the subject of this
thesis. Sums of this type also control the Ky Fan $k$-norms
$\sum_{i\le k}\lambda_i^\ast$ of the singular values, and hence generalise
graph energy~\cite{Nik15}, which explains why the problem has attracted
attention from both the graph-theoretic and the Banach-space sides.

\section{Earlier bounds and projection constants}\label{sec-history}
The first order-linear bounds for individual eigenvalues are due to
Hong~\cite{Hong88}: for $2\le k\le n$,
\begin{equation}\label{hong-bound}
 \lambda_k(G)\le\frac{n-k}2,
\end{equation}
so that $\lambda_2(G)\le n/2-1$, an estimate that is sharp for even $n$ by
$K_{n/2}\sqcup K_{n/2}$. Hong then asked~\cite{Hong93} for the best possible
bounds of this shape. By~\eqref{hong-bound} and the example $K_{n/2}\sqcup K_{n/2}$ one has
$c_2=1/2$. Nikiforov~\cite{Nik15} proved that $c_k$ and $\tau_k$ are finite
for every fixed $k$ and obtained bounds for $c_k$ for all $k>2$, based in part
on Taylor's strongly regular graphs and on Kharaghani's construction of
Hadamard matrices; he also conjectured~\cite[Conjecture 4.2]{Nik15} that for
every $k\ge2$ there is an $\varepsilon_k>0$ with
\begin{equation}\label{nik-conj}
 \tau_k\le\frac{1+\sqrt k}2-\varepsilon_k .
\end{equation}
The reference value $\tfrac12(1+\sqrt k)$ is the one supplied by Mohar's
general estimate~\cite{Mohar}
\begin{equation}\label{mohar-bound}
 M_k(n)\le\frac{1+\sqrt k}{2}\,n ,
\end{equation}
which Mohar proved for graphs and for the wider class of symmetric matrices
with entries in $[0,1]$, and which he asked to be sharpened. For $k=3$ the
reference coefficient is $(1+\sqrt3)/2=1.3660254037\ldots$. The first explicit
improvement below it is the estimate $\tau_3\le3/\sqrt5=1.3416407864\ldots$
proved in Ebrahimi's thesis~\cite[Theorem 3.4.1]{Ebrahimi}. The two-term
matrix analogue, in which the entries are only required to lie in $[0,1]$, was
treated by Ebrahimi, Mohar, Nikiforov and Ahmady~\cite{EMNA}.

The route that currently gives the strongest bounds was opened by Sivashankar,
Tang and Wakhare~\cite{STW}. They reduce the graph-eigenvalue problem to an
extremal problem for orthogonal projections and then invoke the absolute
projection constant $\gamma(r)$, normalised here by the requirement that
\begin{equation}\label{gamma-def}
 \|Q\|_1\le\gamma(r)\,n
 \qquad\text{for every orthogonal projection }Q=Q^T=Q^2,\ \tr Q=r,
\end{equation}
with $\|Q\|_1=\sum_{ij}|q_{ij}|$ the entrywise $\ell_1$ norm. Their two main
estimates are
\begin{equation}\label{stw-individual}
 \lambda_k(G)\le\frac{\gamma(k-1)}{2(k-1)}\,n-1,
 \qquad
 \Lambda_k(G)\le\frac{1+\gamma(k)}2\,n-k .
\end{equation}
The second of these is their Proposition~6.1; together with the bound
$\gamma(r)\le(r+\sqrt{r+2})/(1+\sqrt{r+2})$ of Der\k{e}gowska and
Lewandowska~\cite{DL}, which implies $\gamma(r)<\sqrt r$ for every $r\ge2$,
it settles the conjecture~\eqref{nik-conj}. The first estimate in
\eqref{stw-individual} yields the sharp individual bound
$\lambda_3(G)\le n/3-1$; further recent work on $\lambda_3$ alone is due to
Li~\cite{Li3} and to Leonida and Li~\cite{LeonidaLi}, and lower bounds for
the individual extrema $c_k$ are surveyed in Linz~\cite{Linz}. For $k=3$,
using the classical value $\gamma(3)=(1+\sqrt5)/2$, the second estimate in
\eqref{stw-individual} reads
\begin{equation}\label{stw-hist}
 \Lambda_3(G)\le\frac{3+\sqrt5}4\,n-3=1.309016994374947\ldots\,n-3 .
\end{equation}
For the two-term sum the situation is essentially complete: Kumar, Liu,
Monterde, Pragada and Tait~\cite{KLMPT} proved $\Lambda_2(G)\le8n/7$, Huang
and Wei~\cite{HW} determined the exact maximum together with the unique
extremal graph, whose three part sizes are governed by $n$ modulo $7$, and
Paulin~\cite{Paulin} determined the closures of the rescaled pairs
$\bigl((\lambda_1+1)/n,(\lambda_2+1)/n\bigr)$ and
$\bigl((\lambda_2+1)/n,(\lambda_3+1)/n\bigr)$, from which
$\lambda_1+\lambda_2\le8n/7-2$, $\lambda_3\le n/3-1$ and
$\lambda_2+\lambda_3\le2n/3-2$ all follow. None of these settles
$\Lambda_3$, for the reason explained in the next section.

The graph-limit setting used throughout is standard; we follow
Lov\'asz~\cite{LovaszBook} and Borgs, Chayes, Lov\'asz, S\'os and
Szegedy~\cite{BCLSV}.

\section{Main results}\label{sec-main-results}
% Ported verbatim from paper/top3_joint_moment_bound.tex lines 54-175.
% Only cross-reference wording was adapted; no mathematics was altered.
For a finite simple undirected graph $G$ of order $n\ge3$, put
\[
 \Lambda_3(G)=\lambda_1(G)+\lambda_2(G)+\lambda_3(G),
\]
where the adjacency eigenvalues are ordered algebraically. The signs in
this sum are retained. The extremal problem asks how large $\Lambda_3(G)$
can be when only $n$ is prescribed. The spectral-sum problem appears in Ebrahimi's thesis~\cite{Ebrahimi},
whose Theorem~3.4.1 gives the explicit asymptotic coefficient
$\limsup_n\max_{|G|=n}\Lambda_3(G)/n\le3/\sqrt5<1.3417$;
general bounds for sums of leading eigenvalues were developed by
Mohar~\cite{Mohar}. More recently,
Sivashankar, Tang and Wakhare~\cite[Proposition 6.1]{STW} proved
\begin{equation}\label{stw-bound}
 \Lambda_3(G)\le\frac{3+\sqrt5}{4}n-3.
\end{equation}
Their argument uses the rank-three absolute projection constant
$\gamma(3)=(1+\sqrt5)/2$; a proof of the requisite projection bound is
provided by Der\k{e}gowska and Lewandowska~\cite{DL}.

The two quantities in that reduction are naturally coupled. If
$X\in\R^3$ satisfies $\E XX^T=I_3$ and $X'$ is an independent copy, they are
\[
 A=\E|X\cdot X'|,\qquad B=|\E X|^2.
\]
The separate estimates $A\le(1+\sqrt5)/2$ and $B\le1$ give
\eqref{stw-bound}. A better graph bound therefore requires a quantitative
restriction on their simultaneous size. The mixed projection constant
expressing this question was already introduced in~\cite[Section 6.2]{STW}.
We work directly with its isotropic-law formulation.

Define
\begin{equation}\label{main-constant}
 c=\frac{106106902139519}{107310156800000}
       +\frac{41\sqrt3}{240}-\frac{53}{9000000}
   =1.284673255723918505\ldots.
\end{equation}
\begin{theorem}\label{main}
Every finite simple unweighted graph $G$ on $n\ge3$ vertices satisfies
\begin{equation}\label{main-bound}
 \Lambda_3(G)\le cn-3.
\end{equation}
Moreover, every Borel probability law on $\R^3$ satisfying
$\E XX^T=I_3$ obeys
\begin{equation}\label{law-bound}
 \E|X\cdot X'|+|\E X|^2\le2c.
\end{equation}
Only a finite second moment is assumed; zero vectors and unbounded radii
are allowed.
\end{theorem}
The coefficient decrease from~\eqref{stw-bound} is
$0.024343738651028919\ldots$, with the constant term $-3$ unchanged.
This comparison concerns the stated universal order-only bound. It does
not compare all estimates involving additional graph parameters, nor
assert that $c$ is optimal. For completeness we record the full chain of
explicit coefficients: $3/\sqrt5=1.341640786\ldots$
from~\cite[Theorem~3.4.1]{Ebrahimi}, then $(3+\sqrt5)/4$
from~\cite[Proposition~6.1]{STW}, and finally $c$.

Two adjacent results frame the standing of~\eqref{main-bound}. For the
two-term sum $S_2=\lambda_1+\lambda_2$, Huang and Wei~\cite{HW} determine
the exact maximum together with its unique extremal graph, whose three part
sizes are governed by $n$ modulo $7$; Paulin~\cite{Paulin} determines the
closures of the rescaled pairs
$\bigl((\lambda_1+1)/n,(\lambda_2+1)/n\bigr)$ and
$\bigl((\lambda_2+1)/n,(\lambda_3+1)/n\bigr)$, giving
$\lambda_1+\lambda_2\le8n/7-2$, $\lambda_3\le n/3-1$ and
$\lambda_2+\lambda_3\le2n/3-2$. Neither settles the three-term sum, because
the reduction behind~\eqref{stw-bound} couples two quantities, and it is
their joint control that forces the different geometry analysed here.

One prior announcement bears directly on our lower bound, and we record it
here rather than in a footnote. The self-archived Zenodo
deposits~\cite{Newton1,Newton2} of Newton, dated 4~and 24~August 2026 and
therefore earlier than this work, introduce the same extremal constant
$\cstar$ under the name of a pole--pentagon law, together with the same
optimising weight, and assert the matching upper bound
$\Lambda_3(G)\le\cstar n-3$; the assertion is conditional by its own terms,
since the final positivity layer is a machine certificate whose
admissibility the author explicitly leaves to independent audit
(Remark~\ref{rem-newton}). We claim no priority for the constant.
Theorem~\ref{main} is unconditional and self-contained, and is weaker in
the coefficient by $c-\cstar=0.003105058851750628\ldots$.

We use graphons to pass to an attained maximum. For a symmetric measurable
kernel $W:[0,1]^2\to[0,1]$, let
\begin{equation}\label{sig-definition}
 \Sigma_k(W)=\sup_{\E FF^T=I_k}
       \iint W(x,y)F(x)\cdot F(y)\,dx\,dy,
 \qquad \mathcal M=\sup_W\Sigma_3(W).
\end{equation}
The supremum is over real orthonormal $k$-frames. Since $W\in L^2([0,1]^2)$,
the operator $T_W$ is Hilbert--Schmidt and therefore compact; it is also
self-adjoint, so its spectrum is real and accumulates only at zero, and
$\Sigma_k$ is the sum of the first $k$ positive eigenvalues, padded by zero
if there are fewer than $k$ of them. This convention applies only to
graphons, not to $\Lambda_3(G)$.
The maximum $\mathcal M$ is attained. For a maximizing graphon, choose an
orthonormal top-three eigenframe $F=(f_1,f_2,f_3)^T$, with $f_1\ge0$,
and call $F_\#dx$ its canonical law.

\begin{theorem}\label{structure-main}
Let $W$ attain $\mathcal M$, let $\sum_i p_i\delta_{v_i}$ ($p_i>0$) be its
canonical law, and let $L$ be the diagonal matrix of the three leading
eigenvalues of $W$. Then:
\begin{enumerate}[leftmargin=2.2em,itemsep=2pt]
\item[(a)] the support has at most $24$ points;
\item[(b)] $W$ agrees almost everywhere with the binary step kernel
$\mathbf1_{\{v_i\cdot v_j>0\}}$ on the corresponding positive-measure
fibers;
\item[(c)] the support points satisfy
\begin{gather*}
 \sum_i p_iv_iv_i^T=I_3,\qquad v_i^TLv_i=\mathcal M,\qquad v_{i1}>0,\\
 v_i\cdot v_j\ne0,\qquad
 B_{ij}=\mathbf1_{\{v_i\cdot v_j>0\}},\qquad
 \sum_jp_jB_{ij}v_j=Lv_i;
\end{gather*}
\item[(d)] for every positive-dot clique $J$ in the support, the matrices
$\{v_iv_i^T:i\in J\}$ are linearly independent; in particular $|J|\le6$.
\end{enumerate}
\end{theorem}
This assertion concerns every original maximizer, not merely the
existence of a maximizer with small support.


\section{Method and relation to earlier work}\label{sec-method}
% Ported verbatim from paper/top3_joint_moment_bound.tex lines 177-206.
% Only cross-reference wording was adapted; no mathematics was altered.
The proof has two parts. First, second variation on the actual probability
space gives a negative semidefinite kernel. On each of four sign classes,
this controls the rank of the positive-part Gram kernel and forces finite
canonical support. Moment-preserving variations then imply vector
relations against every binary row of the maximizing kernel. Those rows
separate distinct support points. A minimal-dependency argument shows
that no second-moment relation can survive on a positive clique.

The variation, cloning, and orthant-reduction mechanisms have close
antecedents in Kumar, Mohar, Mojallal and Pragada~\cite{KMMP}. Related graphon
ellipse and clique-cloning arguments for the two-eigenvalue problem
appear in Kumar, Liu, Monterde, Pragada and Tait~\cite{KLMPT}. General
bounds for minimal dependencies of powers of linear forms are due to
Li and Liu~\cite{LiLiu}. Our separating-observation lemma is stated and
proved with its specific hypotheses in Section~\ref{sec-separation};
we do not claim a new general bound for such dependencies.

Second, for one isotropic law we retain the mean deficit, radial variance,
regression error, and weighted harmonic moments of degrees three and four.
A signed quartic majorant couples them in one exact identity. Directional
Hermite inequalities and a radial comparison lead to an explicit sum of
nonnegative terms. Harmonic moment and polynomial-majorant methods are
classical; see Glazyrin~\cite{Glazyrin}, Der\k{e}gowska and
Lewandowska~\cite{DL,DL-SW}, and Sivashankar, Tang and Wakhare~\cite{STW}.
The contribution here is the particular joint control of these errors
for the mixed signed and absolute energy, rather than the use of
polynomial majorants by itself. The support theorem supplies a small
additional term through self-pairs of actual atoms. Diagonal removal and
effective-cardinality estimates are standard weighted-energy devices;
see~\cite{Weighted}.


Chapters~\ref{ch-variational} and~\ref{ch-support} prove the structural
result; Chapters~\ref{ch-moment} and~\ref{ch-bound} give the joint moment
certificate and transfer it to finite graphs. Chapter~\ref{ch-numerical}
reports the verification experiments and the exact scope of the bound, and
Chapter~\ref{sec-further} collects the open problems.

\section{Organization of the thesis}\label{sec-organization}
Chapter~\ref{ch-prelim} collects the standard material used later: the Ky Fan
maximum principle, graphons and the cut metric, equitable partitions and
blow-ups, isotropic laws and absolute projection constants, and the two
harmonic tensor kernels $g_3$ and $g_4$ on the sphere. Nothing in that
chapter is new, and every statement is cited.

Chapter~\ref{ch-variational} sets up the variational quantity
$\mathcal M=\sup_W\Sigma_3(W)$ over graphons, proves that it is attained,
gives the two-sided estimate
$\cstar\le\mathcal M\le(1+\sqrt3)/2$, and derives the four structural facts
that every maximiser satisfies: the ellipsoid identity $F^TLF=\mathcal M$,
the row-insertion inequality, the threshold form of the kernel, and the
simplicity relations $\lambda_1>\lambda_2>\lambda_3>\lambda_4$.

Chapter~\ref{ch-support} proves that the canonical law of every maximiser is
finitely supported, with at most $24$ atoms. The argument runs through a
second-variation inequality on the actual support, propagation of moment
relations, a dimension descent over sign regions, and a separating family of
binary observations that excludes minimal second-moment dependencies on
positive cliques.

Chapter~\ref{ch-moment} builds the moment machinery for a single isotropic
law on $\R^3$: the harmonic decompositions, the regression estimate, the
signed quartic majorant $f$, and the exact joint defect identity that ties the
mean deficit, the radial variance, the regression error and the two harmonic
moments together.

Chapter~\ref{ch-bound} converts that identity into a coefficient bound.
Directional Hermite inequalities and a radial comparison produce an explicit
sum of nonnegative terms; a second pass extracts a reserve from the quartic
slack; the support theorem contributes a positive atomic term; and the final
certificate gives $\mathcal M\le c$. The chapter closes with the transfer to
all isotropic laws and to finite graphs, the proof of
Theorem~\ref{main}, and the weighted-graph form of the bound.

Chapter~\ref{ch-numerical} is computational. It records an exact-arithmetic
verification protocol, an exhaustive check over all graphs of orders
$3\le n\le7$, an exact spectral formula for the equal blow-ups of
$K_1\vee C_5$, convergence tables for the optimally weighted blow-ups, and a
discussion of what these experiments do and do not establish.

Chapter~\ref{sec-further} states the conclusions and lists the open problems,
including the residual gap $c-\cstar$, the classification of extremal
templates, and the adjacent two-term problems. Appendix~\ref{sec-appendix}
records which constants in the argument are forced and which were chosen and
then verified.

\section{Notation and conventions}\label{sec-notation}
A complete index is given in the List of Notation after the table of
contents; the few conventions that matter globally are these. All graphs are
finite, simple and undirected, and $n$ always denotes $|V(G)|$ with $n\ge3$.
Eigenvalues are ordered algebraically, not by absolute value, and
$\Lambda_k$ retains signs. The letter $A$ denotes an adjacency matrix and
$T_W$ the integral operator of a graphon $W$. For a real symmetric matrix
$Q$, $\|Q\|_1=\sum_{ij}|Q_{ij}|$ is the entrywise $\ell_1$ norm and
$\|Q\|_F^2=\tr(Q^2)$ is the squared Frobenius norm; these should not be
confused with the operator norms $\|\cdot\|_{2\to2}$ and $\|\cdot\|_\square$.
The symbol $\E$ denotes expectation, and $X'$ is always an independent copy
of $X$. Rational numbers written with large numerators and denominators are
exact, not rounded; rounded values are always followed by an ellipsis.

Four labels are used to distinguish the epistemic status of a statement.
\emph{Proved} means a complete argument appears in this thesis.
\emph{Verified} means the statement is an explicit identity or inequality
that was re-checked in exact arithmetic. \emph{Numerical} means floating
point evidence only. \emph{Gap} means an assertion that is neither proved nor
verified here, and every gap is repeated in Chapter~\ref{sec-further}.
